\chapter{}
Poucos meses depois do início das aulas no São Francisco de Sales, ou Diocesano como era conhecido, porque pertencera à Diocese antes de ser repassado aos jesuítas, fui chamada pelo seu Diretor, Padre Ângelo, para substituir um professor de História que abandonara suas turmas sem maiores explicações.
Aceitei, satisfeita.
Sentia falta do meu trabalho.

Quando comecei a lecionar no Diocesano, faltavam ao velho Colégio menos de duas décadas para completar cem anos.

\begin{figure}
\centering
\includegraphics[width=\linewidth]{35/fachada.png}
\caption*{A histórica fachada do São Francisco de Sales em Teresina.}
\end{figure}

Foi durante os anos 50, no século passado, que um grupo de jesuítas italianos chegou a Teresina em obediência à missão da Companhia de evangelizar através da educação.
Como a Diocese vinha encontrando dificuldades para manter em funcionamento o São Francisco de Sales, os missionários se prontificaram a assumi-lo.
Padre Ângelo era um remanescente desse grupo.
Havia ainda outros que eu logo conheceria: Padre Florêncio.
Padre Gino e Irmão Guido.
Os outros já tinham morrido, ou tinham sido transferidos para outras obras.

\begin{figure}
\centering
\includegraphics[width=0.9\linewidth]{35/meninos+padres.png}
\caption*{Pe. Ângelo e Pe. Gino com os meninos, na primeira comunhão.}
\end{figure}

Assim que assumi as turmas, fui advertida pelo Padre Diretor de que não precisava me preocupar com inovações pedagógicas.
O método preferencial na escola era mesmo a preleção pura e simples, num ambiente de disciplina e atenção concentrada.
Achei melhor não contestar, mas ignorei solenemente a recomendação.
Juntei meus alunos, levei-os para o pátio, organizei-os em grupos de discussão, ensinei-os a pesquisar, fui procurar pela biblioteca, enfim, estava morrendo de saudades de exercer meu ofício e não ia abrir mão de fazê-lo à minha maneira.
Percebi Padre Ângelo me observando de longe, parecendo intrigado, ou será que afrontado com o meu atrevimento? Fui em frente.
Os alunos pareciam entusiasmados, recebiam bem as novidades, até porque o professor anterior tinha deixado tão péssima lembrança que quem quer que viesse depois pareceria inevitavelmente muito melhor.

Lecionei por pouco tempo e, nesse período, cheguei a dar aulas ao Tomaz e ao Marcelo.
Nunca soube o que Tomaz, sempre discreto, pensava do meu trabalho.
Porém, Marcelo gostou muito da oportunidade de exercitar seu nascente senso crítico.
Marcava cerrado.
Essa experiência se repetiria, anos depois, com Vicente, também um crítico atento.
No fundo, achei formidável que eles se sentissem à vontade para fazê-lo, distinguindo a mãe da professora e, com isso, incentivando os colegas a exercer também seu direito de contestação.
Que mais podia desejar alguém que lecionava História? 

O que rapidamente ficou claro é que eu não podia interferir na avaliação ou nos exercícios prescritos pelos professores, sem discutir com eles o processo de ensino-aprendizagem como um todo.
Continuei a dar algumas aulas, mas, de professora de História fui me transmutando paulatinamente em coordenadora pedagógica, metamorfose só viabilizada porque os jesuítas acreditavam na formação em serviço.
Esse trabalho, que eu exerceria pelo resto da minha vida profissional, foi altamente gratificante porque, além de me permitir pôr em prática medidas que eu acreditava imprescindíveis a um trabalho atualizado e responsável em Educação, tornou possível realizar meu sonho de garantir aos meus filhos e ao restante dos alunos o acesso a um ensino superior de qualidade.
Mas, para isso, um longo caminho foi trilhado.

A Companhia de Jesus no Brasil, e especialmente no Nordeste, começava a incentivar um esforço de modernização da prática pedagógica nas suas escolas.
Isso significava dizer que a pedagogia inaciana, um conjunto de práticas nascidas da formidável intuição de Inácio de Loyola, o fundador da ordem, precisava ser revisitada e reexaminada á luz dos acontecimentos que incidiram sobre o mundo na virada do século.
O momento era propício, ademais, porque os jesuítas europeus, guardiães ortodoxos da visão inaciana de educação, já muito idosos, começavam a ser substituídos na direção dos colégios pelos religiosos formados no Brasil.
Este era o caso do São Francisco de Sales, onde Padre Ângelo foi, pouco tempo depois, substituído por Padre Darly, um capixaba que trouxe como auxiliar outro brasileiro e piauiense, o Irmão Pires, psicólogo com formação em Teologia e Filosofia, cuja orientação foi fundamental tanto para o meu crescimento pessoal, quanto profissional.
 
A progressiva retirada dos jesuítas europeus, contudo, ameaçava abrir uma grave lacuna na educação ministrada nas escolas da Companhia, pelo fato de que, na sua ortodoxia, eram exigentes em disciplina e método, dois princípios muito necessários à formação de indivíduos autônomos e eficientes.
Nós, os brasileiros, costumamos creditar nossas graves deficiências de uma e de outro a uma pretensamente encantadora ``informalidade''.
Disso resulta que nossas cabeças pensantes, à parte honrosas exceções, são, a torto e a direito, facilmente seduzidas por juízos grandiloquentes, teorias radicais, afirmações bombásticas sem nenhuma comprovação e pior, agem a partir disso, não vendo no fato qualquer inconveniência.
Sem falar na incapacidade endêmica de planejar e levar a cabo qualquer empreitada sem as costumeiras idas e vindas, correções de rota, dilatações de prazos e exigências infindáveis de recursos, isso quando não se larga tudo pelo meio e pronto.
Felizmente, para o nosso Colégio, nem todos os italianos se foram.
Um ficou, o mais irredutível deles, a cuja influência devo uma parte preciosa do caráter dos meus filhos: Padre Florêncio.

\begin{figure}
\centering
\includegraphics[width=\linewidth]{35/paula+marcelo+padre.png}
\caption*{Paula e Marcelo visitando Pe. Florêncio.}
\end{figure}

Um dia, estava eu trabalhando na minha sala, quando o velho jesuíta irrompeu porta adentro, os olhos azuis faiscando, a longa barba tremendo de uma fúria contida a custo.
Até então, só nos cumprimentávamos pelos corredores.
Atirou sobre a mesa um calhamaço de provas e vociferou:

\textit{``-- Veja a senhora se é possível! Esses meninos tomam diariamente 250 mililitros de refrigerante na cantina e, apesar disso, acham que a sala onde assistem aula tem 150, ou 75 ou 300 mililitros de volume! O que lhes foi ensinado de Matemática até agora?''}

Tratava-se de um primeiro colegial em que ele havia aplicado a avaliação em Sistemas de Medidas, que infalivelmente precedia o curso de Química ministrado por ele.
Padre Florêncio regia também o laboratório da disciplina, ao qual só tinham acesso os alunos escolhidos por ele.

Uma rápida olhada nas respostas dos alunos deixou claro que as mudanças na escola tinham que começar imediatamente.
Entre a poderosa inteligência do Irmão Pires, seu profundo domínio do pensamento inaciano, sua perspicaz compreensão dos homens e a incorruptível coerência de Padre Florêncio e sua honesta devoção à missão de ensinar e contando, ainda, com a ajuda de um punhado de corajosos educadores piauienses, começamos o trabalho de modernizar o Colégio.

De um cuidado eu não abria mão: o de preservar, no que tinha de essencial, a visão inaciana de educação.
Sobretudo porque conhecê-la deu forma a intuições que me orientavam na educação dos meus filhos e dos meus alunos.
O objetivo precípuo do desenvolvimento da autonomia, a valorização da emoção como garantia de adesão ao compromisso, o exercício constante e obrigatório do discernimento como preparação para ação, tudo isso, descobri, era uma fantástica maneira de ajudar um indivíduo a crescer ético, íntegro e equilibrado.

Nas minhas leituras sobre sabedoria rabínica, narra-se que o Faraó teria indagado de Moisés o que ele seria capaz de fazer para surpreender a si próprio, uma vez que para surpreender os outros basta fazer aquilo de que não nos julgam capazes, mas surpreender-nos a nós próprios é muito mais difícil, porque significa ir além do que acreditávamos saber ou poder.
Se essa pergunta fosse feita a mim, eu talvez me referisse ao que realizei naqueles anos de trabalho no Colégio.

Meu temperamento caracteristicamente introvertido sempre me fez avessa a qualquer forma de envolvimento social.
Curiosamente, por um total de catorze anos em que trabalhei no Diocesano, eu viveria cercada de pessoas, eu lideraria pessoas, eu me envolveria com um grande universo de pessoas e faria isso com enorme entusiasmo e com pleno sentimento de realização.
Porque a verdade é que, à parte o fato de ter me tranquilizado quanto à qualidade do estudo dos meus filhos, aquele período ajudou-me a encontrar minha identidade pela necessidade que eu tive de conferir maior nitidez e fundamentação às minhas convicções, ao mesmo tempo em que me tornou mais consciente das minhas limitações e das minhas possibilidades.
Ao longo de todos aqueles anos em que levamos adiante o projeto de atualização da pedagogia inaciana li muito, discuti muito, conheci pessoas admiráveis e fui desenvolvendo ações e assumindo posturas que alicerçaram meu desenvolvimento pessoal e me ajudaram a lidar melhor com a vida, com o mundo e com os meus semelhantes em geral.

Hoje, defendo que todas as mulheres devem buscar um espaço de realização pessoal, não importa se dentro ou fora de casa.
Não se trata, a meu ver, de se igualar aos homens nos direitos e regalias, de vencer preconceitos, mostrar-se capaz, eficiente, competitiva e sei lá o que mais.
Isto me parece absolutamente secundário.
Mesmo porque, o poder de gestar Vida que lhe pertence com exclusividade, faz dela um ser privilegiado por natureza e qualquer mulher que tenha clara essa consciência sabe do que estou falando.
Entretanto, até para desfrutar em plenitude esse privilégio, é preciso que ela vá ao encontro do mundo na busca de conquistar sua integridade, descobrir sua unicidade.
Homens e mulheres, indistintamente, precisam encontrar seu lugar no grande plano da Criação, o lugar que cada um deles precisa ocupar individualmente, exclusivamente.
Cada um alinhado à sua finalidade, para se sentirem realizados na sua humanidade, como ensinou Santo Tomás de Aquino.

E o mais formidável nessa quadra da minha vida é que, ao contrário do que eu temia, meus filhos se comportaram da forma mais compreensiva e cooperativa que se possa imaginar.
Sem exceção, mantiveram um desempenho escolar irrepreensível, malgrado os eventuais aborrecimentos que tiveram por serem ``os filhos da Coordenadora'', o que deu força ao meu trabalho, pois foram a prova viva de que as exigências que eu fazia aos demais eram perfeitamente realizáveis por alunos normais, capazes de estudar, divertir-se, praticar esportes e cumprir seus deveres, mesmo tendo a mãe trabalhando fora.